\chapter*{Preface}
\addcontentsline{toc}{chapter}{Preface}

This is not an introductory LaTeX manual. Introductions to LaTeX are plentiful, well written, and not the gap this book is trying to fill. Most of them teach a reader how to produce a correct document: how to start a section, insert a table, number an equation, cite a source. That is necessary knowledge, and this book assumes a reader already has it, at least in outline. What that kind of introduction does not usually teach, because it is a different kind of knowledge entirely, is how to design and maintain a LaTeX project that has grown past the scale a single file, a weekend, or a single author's unaided memory can comfortably hold.

This book exists for that second kind of knowledge. Its central claim, argued from several different directions across twenty-two chapters, is that LaTeX is best treated not as a typesetting language to be learned once and then simply used, but as a programmable, reproducible publishing environment, closer in spirit to a software project than to a word-processed document. Once that reframing is taken seriously, a great deal of what otherwise looks like idiosyncratic LaTeX trivia, why a build needs two passes, why a fragile command breaks only sometimes, why a five-hundred-page manuscript should never live in one file, turns out to be the direct, predictable consequence of a small number of underlying principles, the same principles a software engineer would recognize from any other large, long-lived, collaboratively maintained system.

The book is organized to build that understanding from the ground up: how TeX actually processes a document before any LaTeX convention is layered on top of it; how to architect, version, and build a large multi-file project; how to write macros and interfaces that behave predictably; how to use modern font technology and mathematical typesetting to their full extent; how to manage bibliographies, indexes, and glossaries at real scale; how to bring graphics and computation directly into the document's own build process rather than pasting in artifacts prepared elsewhere; and, in its final chapters, how to think about an entire body of published work, spanning years and multiple books, as one maintained, coherent system rather than a shelf of unrelated documents.

Two companion volumes sit alongside this one without being required reading for it. \emph{Thinking in Vim} is about editing structured text efficiently, treating an editor as more than a place to type. \emph{Thinking Like Rust} is about building computation reliably, with the kind of explicit, deterministic discipline that a document's own correctness can depend on when a figure or a table is generated rather than hand-copied. Each book stands on its own, and nothing in this one assumes the other two have been read. But a reader who has encountered all three will recognize the same underlying commitment running through each: that a technical practice, whether editing, computing, or publishing, rewards being learned as a designed system rather than picked up as a bag of separately memorized tricks. This book is the last of the three stages, the one responsible for taking whatever was written and computed and turning it into something durable enough to be read.

\chapter*{Acknowledgements}
\addcontentsline{toc}{chapter}{Acknowledgements}

This book, like the rest of the corpus it belongs to, was developed through iterative drafting and revision rather than in a single sustained sitting, and it owes whatever clarity it has to that process of repeated reworking rather than to any single moment of insight. Readers of earlier drafts of the material collected here, and collaborators on the broader research program this book's conventions were drawn from, have the author's thanks, even where they are not named individually in this edition.

\chapter*{Who This Book Is For}
\addcontentsline{toc}{chapter}{Who This Book Is For}

This book is for a reader who already knows how to write a LaTeX document and now wants to know how to build one that lasts. That includes, without being limited to: an academic or independent researcher maintaining a long-running body of technical writing across many documents; a graduate student or postdoc whose thesis or first book-length manuscript has outgrown the single-file habits that worked fine for a conference paper; a software engineer who has picked up LaTeX for documentation or papers and wants to apply the same engineering discipline to it that they would apply to a codebase; and anyone maintaining a technical book, textbook, or reference work expected to remain buildable, correct, and internally consistent for years after it is first written.

It is not primarily for a reader who has never used LaTeX before, though such a reader is welcome to use this book as a map of where their skills might eventually lead, alongside a more conventional introduction covering the syntax this book largely assumes. It is also not a reference manual for any single package; each chapter points to the specific tools relevant to its subject, but the authoritative documentation for any given package remains the right place to look up that package's full interface. What this book offers instead is the architecture around those tools: how to choose among them, how to combine them into a working system, and how to keep that system working as it, and the body of work built with it, continues to grow.
